\section{Provisioning}

\subsection{Hosting provider}

We chose to rent a VPS from Hetzner due to its strong balance between affordability and performance. Hetzner is well known provider and it's built-in firewall was another key factor. It lets us apply firewall rules at the infrastructure level, filtering unwanted traffic before it reaches the server. After completing the migration and application setup, we updated our Go provisioning script to automatically create and attach the firewall to each server.

\subsection{VPS}

Converting the legacy system from Python to Go significantly improved performance, with faster request handling, lower CPU usage, and reduced memory consumption. This allowed the application to run efficiently on less powerful hardware while handling a similar workload.

We chose a Hetzner CX23 server with 2 vCPUs, 4 GB RAM, and 40~GB storage. Multiple vCPUs were necessary to support concurrent applications, while 4 GB RAM ensured stable performance without heavy swap usage. The server also includes 20 TB monthly egress, offering strong cost efficiency at just \euro3,99 per month.

\subsection{Provisioning the server}

We chose to define our infrastructure provisioning logic as code rather than relying on Hetzner's GUI. Hetzner provides an official Go library that acts as a wrapper around their API, which we used to automate the provisioning of the server directly from our application code.

Terraform was also considered as a potential solution for infrastructure provisioning. However, we ultimately preferred the simplicity and maintainability of using a single Go script. This decision aligned naturally with the broader migration of the main application from Python to Go, allowing the team to remain within a familiar ecosystem and programming language.

\subsection{Bootstrapping}

Server bootstrapping is a critical part of deployment, as misconfigurations can cause downtime, unstable services, and security vulnerabilities. To reduce these risks, we used Ansible for automated server provisioning and configuration.

Ansible's idempotent design lets playbooks run repeatedly without drift, keeping the setup reproducible and reducing manual error.

The setup consists of two Ansible playbooks. The first handles server bootstrapping and security hardening, including system configuration, security measures, and deployment preparation.

\begin{table}[h]
\centering
\begin{tabular}{|p{4cm}|p{9cm}|}
\hline
\textbf{Security Measure} & \textbf{Description} \\
\hline
Non-root administrative user & Creates a user without direct root access while granting \texttt{sudo} privileges. \\
\hline
System upgrades & Updates and upgrades server packages and system components. \\
\hline
Disable root login & Prevents direct root login to improve security. \\
\hline
System hardening & Applies Linux kernel, SSH, and network hardening settings. \\
\hline
UFW and Fail2Ban & Configures firewall rules and intrusion prevention mechanisms. \\
\hline
Auditd setup & Enables system auditing and security event logging. \\
\hline
Unattended upgrades & Configures automatic installation of security updates. \\
\hline
\end{tabular}

\caption{Security hardening measures applied during server bootstrapping}
\label{tab:security_measures}

\end{table}

The second playbook focuses specifically on installing and configuring Docker along with the Docker Compose plugin. Separating these responsibilities into distinct playbooks improves maintainability and makes the provisioning process easier to understand and manage.